\section{Additional experiment results}
\label{sec:additional_quant}
\subsection{Novel view synthesis via latent decoding}
\begin{figure}[t]
    \centering
    \includegraphics[width=\linewidth]{supple_sec/supple_figure/assets/sup_cherry_rae.pdf}
    \vspace{-15pt}
    \caption{\textbf{Novel view synthesis via latent decoding.} We explore the potential of combining our view-invariant feature decoder, \oursF, with generative models. We lift DINOv2-base features~(which serve as latents for a Representation Autoencoder (RAE)~\cite{zheng2025diffusion}) extracted from the input views~((a)) to 3D Gaussians and render them at novel viewpoints. (b) \textbf{Ours -- 3DGS Rendering}: Standard RGB rendering from the estimated Gaussians provides faithful reconstruction but exhibits some blurriness in high-frequency regions. (c) \textbf{Ours -- Latent Decoding}: By decoding the rendered feature maps through the RAE's diffusion transformer with a 1-step denoising process, we recover significantly sharper textures and edges. Although the generative nature introduces slight variations in fine details compared to the Ground Truth (d), this validates our model's ability to provide consistent 3D-aware latents for diffusion-based pipelines.} 
    \label{supfig:rae}\vspace{-10pt}
\end{figure}

As our multi-view invariant feature decoder \oursF\ can take any desired feature as input, we conduct an additional experiment of leveraging the latents of the representation autoencoder~(RAE)~\cite{zheng2025diffusion} as input. Specifically, we leverage the open-sourced RAE, which learns a diffusion transformer~\cite{peebles2023scalable} and a decoder based on the frozen DINOv2-base~\cite{oquab2023dinov2} as the latent. By lifting the DINOv2-base features and rendering novel viewpoints, we can decode the rendered features using the pre-trained RAE decoder or DiT after adding a small noise to the rendered features. In Fig.~\ref{supfig:rae}, we show the results of novel view synthesis obtained from the estimated 3D Gaussians~(Fig.~\ref{supfig:rae}-(b)) and the novel view synthesis obtained by decoding the rendered features through the DiT~(Fig.~\ref{supfig:rae}-(c)). Specifically, for this experiment, we add a 1-step noise from the original 50-step denoising schedule of RAE-DiT~\cite{zheng2025diffusion}, and denoise with \mbox{1-step}. As shown in Fig.~\ref{supfig:rae}, renderings from 3D Gaussians already provide sufficiently good results but still contain some blurry regions, whereas the latent decoding yields significantly sharper results. However, since the pipeline involves a generative process, even with a single step, we observe slight deformations in fine details. Nevertheless, this experiment demonstrates the potential of combining novel view synthesis with diffusion processes by leveraging latents as input to our \oursF.

\subsection{Novel view synthesis in two-view setting}
In this section, we evaluate novel view synthesis performance on the RealEstate10K dataset~\cite{zhou2018stereo}, following NoPoSplat's~\cite{ye2024no} protocol where two input images are used to estimate 3D Gaussians and render a target view. We compare with both \textit{pose-dependent} models~\cite{charatan2024pixelsplat, chen2024mvsplat} and \textit{pose-free} models~\cite{hong2024pf3plat,hong2024unifying, smart2024splatt3r, ye2024no,huang2025no}. Since we use VGGT~\cite{wang2025vggt} as our visual encoder for \ours, we additionally train VGGT+NoPo, which replaces NoPoSplat's MASt3R~\cite{leroy2024grounding} backbone with VGGT~\cite{wang2025vggt} while maintaining NoPoSplat's pipeline to estimate per-pixel Gaussians. Note that \ours\ directly estimates Gaussians from unposed images, falling into the \textit{pose-free} category. For NoPoSplat, VGGT+NoPo, and ours, we follow NoPoSplat's test-time camera pose optimization, which is only necessary for evaluation. As shown in Tab.~\ref{tab:nvs_baseline}, despite estimating \textbf{65$\times$} fewer Gaussians than per-pixel methods~\cite{hong2024pf3plat,hong2024unifying, smart2024splatt3r, ye2024no,huang2025no}, our approach achieves comparable rendering quality with much faster speeds, validating that our compact Gaussians is sufficient for 3D scene reconstruction.

\begin{table}[t]
\centering
\caption{\textbf{Comparison of novel view synthesis on RealEstate10K~\cite{zhou2018stereo}.} Our method maintains competitive results while using far fewer Gaussians.}
\label{tab:nvs_baseline}
\resizebox{\linewidth}{!}{
\begin{tabular}{c|l|cccccc}
\toprule
\multirow{2}{*}{\shortstack{{Pose-}\\{free}}} & \multirow{2}{*}{Methods} & \multicolumn{6}{c}{Average} \\
& & \#G $\downarrow$ & Memories$\downarrow$ & FPS$\uparrow$ & PSNR$\uparrow$ & SSIM$\uparrow$ & LPIPS$\downarrow$ \\
\midrule
\multirow{2}{*}{\ding{55}} & PixelSplat~\cite{charatan2024pixelsplat} & 131K & 33.6MB &388.03 &  23.848  & 0.806 & 0.185 \\
& MVSplat~\cite{chen2024mvsplat}& 131K & 33.6MB &392.6 & 23.977 & 0.811 & 0.176 \\
\midrule
\multirow{9}{*}{\ding{51}} & CoPoNeRF~\cite{hong2024unifying}  & -& -&0.4 & 18.938 & 0.619 & 0.388\\
& Splatt3R~\cite{smart2024splatt3r} &131K & 33.6MB &393.1& 15.318 & 0.490 & 0.436 \\
& PF3plat~\cite{hong2024pf3plat}&131K & 33.6MB &397.1 & 21.042 & 0.739 & 0.233 \\
& SPFSplat~\cite{huang2025no}&131K & 33.6MB &397.3 & 25.845& 0.852& 0.152\\
& NoPoSplat~\cite{ye2024no} & 131K & 33.6MB &369.8& 25.033& 0.838& 0.160\\
& VGGT+NoPo~\cite{wang2025vggt, ye2024no}& 100K & 25.6MB &419.8 & 23.015& 0.762& 0.187\\
& AnySplat~\cite{jiang2025anysplat} & 320K & 81.9MB & 319.2 & 18.828& 0.656 & 0.358\\
& \textbf{C3G (Ours)} & 2K & 0.1MB & 451.7& 22.387& 0.713& 0.259\\
\bottomrule
\end{tabular}}
\vspace{-10pt}
\end{table}